\section{Lumped elements at a transverse cut}
\label{sec:loads}

\noindent\emph{Status: draft -- two element kinds are reported here, a series resistance and a shunt
capacitance, both placed at a transverse cut across the strip of Section~\ref{sec:msres-box}. Every series
figure is evaluated against a load-free run of the same line, and every shunt figure against the capacitance the
mesh was asked to inject. The series branch passes the test it was registered against. The shunt branch does not
pass, and is reported as an envelope that says how far it does get and no further.}

\subsection{What a cut means numerically}
\label{sec:loads-cut}

A \texttt{series\_r} load of value \texttt{inf} is a bare open, and the way it is built has to be stated before
any number is quoted. Nothing is assembled at the load position: the cut annihilates the perfectly conducting
mask on one column of dual edges, so the strip stops there and the two halves of the line are joined by nothing at
all. This is not a simplification of a component the solver cannot discretise. The cut \emph{is} the model, and
the only thing standing between the two halves of the strip is the field of the cut itself.

That field is not negligible. A column of edges whose conductor has been removed carries a capacitance of its own
-- the capacitance of the gap across the strip -- and it carries it whether or not a user asked for one. The bare
open therefore does not present a pure open circuit to the line: it presents that gap capacitance in parallel
with whatever the load was meant to be, and a prediction that treats the cut as an ideal series element is wrong
by that amount. What follows measures the term and finds where it stops being a small correction.

\subsection{The series branch}
\label{sec:loads-series}

Three resistances are cut into the strip at $15.20$ mm and each is measured against a load-free run of the same
line, so that the reference branch carries the cut's capacitance and the load branch differs from it in the load
alone. The board and both of its reference numbers are named in the caption of
Table~\ref{tab:loads-series}, and every figure below is read from that record rather than re-derived.

\begin{table}[H]
  \centering
  \caption{The series branch on the reference board: $0.400$ mm cells, $n_{\mathrm{sub}} = 4$,
  $n_{\mathrm{strip}} = 8$, the cut at $15.20$ mm, $2400$ steps, each resistance evaluated against a load-free
  run of the same line. That run drives at $5.8278$ GHz and gives a station impedance
  $Z_c = 50.547 + 4.025j\ \Ohm$, and the cut's capacitance from the mesh formula
  $C_g = n_{\mathrm{strip}}\eps_0\eps_r\,\dd x = 1.2467\times10^{-13}$ F $= 0.1247$ pF, so that
  $Z_g = 1/(j\omega C_g) = -219.061j\ \Ohm$. The fourth column is the textbook series formula evaluated with that
  $Z_c$; the last is the deviation of the measurement from it. The pre-registered prediction for the
  $100\ \Ohm$ point was a deviation of $7.1$ per cent.}
  \label{tab:loads-series}
  \begin{tabular}{@{}rrrrr@{}}
    \toprule
    $R$ [$\Ohm$] & $|s_{21}|$ measured & $\arg s_{21}$ [$^\circ$] & $2Z_c/(2Z_c+R)$ & deviation [pc] \\
    \midrule
     25 & 0.794825 & $+1.4852$  & 0.802640 & $-0.974$ \\
     50 & 0.673329 & $+4.9999$  & 0.670249 & $+0.460$ \\
    100 & 0.549464 & $+13.0870$ & 0.503909 & $+9.040$ \\
    \bottomrule
  \end{tabular}
\end{table}

The textbook formula holds to under one per cent at $25$ and $50\ \Ohm$, and it fails the $100\ \Ohm$ point by
$9.04$ per cent. That failure was not a surprise found afterwards: the discriminator was registered at three
resistances in advance, and the prediction for the third was a few per cent of deviation from the ideal series
reading. The reason the formula fails is the one Section~\ref{sec:loads-cut} gave: it treats the cut as an ideal
series impedance, whereas the one-cell cut also carries $C_g$, so the load the line actually meets is $R$ in
parallel with $Z_g$.

That reading is testable without appealing to the formula under test, by inverting the measured transmission
coefficient for the impedance the line saw, $Z_l = 2Z_c(1/s_{21} - 1)$, and comparing that impedance with
$R \parallel Z_g$ computed from the geometry alone.

\begin{table}[H]
  \centering
  \caption{The same board and the same three resistances, with the load impedance inverted from the measured
  transmission coefficient and compared with $R$ in parallel with the cut capacitance $Z_g = -219.061j\ \Ohm$ of
  Table~\ref{tab:loads-series}. The columns are, left to right, the commanded $R$, the inverted $Z_l$,
  $R \parallel Z_g$, their difference, its modulus, and the deviation of the measurement from
  $2Z_c/(2Z_c + R \parallel Z_g)$; the largest phase difference against that corrected prediction, over the
  three points, is $0.68^\circ$, and the corrected values themselves are $0.805632$, $0.682260$ and $0.544544$.}
  \label{tab:loads-implied}
  \small
  \begin{tabular}{@{}rrrrrr@{}}
    \toprule
    $R$ [$\Ohm$] & $Z_l$ [$\Ohm$] & $R \parallel Z_g$ [$\Ohm$] & diff. [$\Ohm$] & $|$diff.$|$ [$\Ohm$] & corr. dev. [pc] \\
    \midrule
     25 & $26.3162 - 1.2222j$ & $24.6786 - 2.8164j$ & $+1.6377 + 1.5942j$ & 2.29 & $-1.34$ \\
     50 & $49.5174 - 9.2259j$ & $47.5242 - 10.8472j$ & $+1.9932 + 1.6214j$ & 2.57 & $-1.31$ \\
    100 & $81.4314 - 35.4408j$ & $82.7550 - 37.7771j$ & $-1.3235 + 2.3363j$ & 2.69 & $+0.90$ \\
    \bottomrule
  \end{tabular}
\end{table}

The corrected prediction is within $1.34$ per cent at $25\ \Ohm$, $1.31$ per cent at $50\ \Ohm$ and $0.90$ per
cent at $100\ \Ohm$, against deviations of $0.97$, $0.46$ and $9.04$ per cent for the ideal series formula, and
every phase agrees to within $0.68$ degrees. What is left over is a nearly constant offset rather than a
residual that grows with the load: the real part of the difference in Table~\ref{tab:loads-implied} runs from
$-1.32$ to $+1.99\ \Ohm$ and its imaginary part from $+1.59$ to $+2.34\ \Ohm$ across the three, which is the
signature of a small fixed series term in the cut, not of an error in the capacitance. The capacitance the mesh
formula predicts is therefore confirmed by measurement, and the confirmation is worth as much as the test it was
registered against: at $50\ \Ohm$ the load and the line impedance are nearly equal, so the ideal and corrected
formulas agree there almost exactly, and \emph{a single $50\ \Ohm$ probe would have appeared to confirm the
textbook formula by cancellation}. That is why the discriminator was registered across three resistances and
not one.

One detail of the reference branch is a choice and is declared as one. It takes the attenuation for its
propagation constant from the joint complex-$\gamma$ fit of the station series, $3.0436$ Np/m, and the
normalisation of the ratio is built so that the choice of that attenuation is immaterial to what the ratio
reports.

\subsection{What the reference divides out}
\label{sec:loads-ratio}

The series resistances above were each measured against a load-free run, and the size of that correction is easy
to underestimate: a loaded record cannot be normalised against itself. On the record named in the caption of
Table~\ref{tab:loads-ratio} one loaded run is quoted three ways: with the load-free twin supplied as its reference,
with no reference at all, and as the twin itself on the identical mesh. The three readings below are one record, not
three runs of three line lengths.

\begin{table}[H]
  \centering
  \caption{The reference measurement on a microstrip record: $\eps_r = 4.4$, $h = 1.60$ mm,
  $n_{\mathrm{strip}} = 8$ cells across the strip, six matched-layer cells on every face, $300$ steps, the drive
  snapped to $5.827786$ GHz, and a $100\ \Ohm$ series resistance at $15.20$ mm along a $30$ mm line. All three
  rows are the same record: the first is the load-free twin on the identical mesh, the second is the loaded run
  handed that twin as its reference, the third is the same loaded run with no reference given. All three are
  $300$-step readings of one record, and what they demonstrate is that the normalisation is self-consistent, not
  what the load transmits: record length is not converged at $300$ steps, and Table~\ref{tab:loads-length}
  measures how far from converged the readings are.}
  \label{tab:loads-ratio}
  \begin{tabular}{@{}lr@{}}
    \toprule
    reading of the same record & $|s_{21}|$ between the station planes \\
    \midrule
    load-free twin             & $0.541741$ \\
    loaded, reference given    & $0.613960$ \\
    loaded, no reference given & $0.327195$ \\
    \bottomrule
  \end{tabular}
\end{table}

The record length is the one variable that separates these readings from the ones quoted for the load, and it has
to be measured rather than assumed away. The mesh is fixed in every row of Table~\ref{tab:loads-length} -- the
same $0.400$ mm cell, the same six matched-layer cells on every face, the same $100\ \Ohm$ series element at
$15.20$ mm -- and what changes is $n_{\mathrm{steps}}$, of which the transform uses $0.75\,n_{\mathrm{steps}}$
samples. Between $300$ and $9600$ steps the loaded reading moves from $0.613960$ to $0.556575$, while the
load-free twin moves from $0.541741$ to $0.989576$. From $2400$ steps upward the loaded series is monotonically
increasing, with a step that shrinks from $1.128$ to $0.164$ per cent per doubling of the record length, which
puts its limit near $0.5567$ and puts the $2400$-step figure about $1.3$ per cent below that limit. The twin's
step shrinks more slowly, $0.650$ to $0.196$ per cent, and at $9600$ steps it is still creeping up, so the twin is
not converged either. The consequence has to be stated plainly: the denominator of the ratio is itself a function
of record length, so the division cancels a factor that is $O(1)$ different at $300$ steps, and the identity
quoted below therefore holds while both of its terms are still unsettled. For this element the value to quote is
the converged loaded reading, about $0.5567$; the $0.549464$ of Table~\ref{tab:loads-series} is a $2400$-step
reading of the same element, $1.3$ per cent below that limit, and it is not a converged number.

\begin{table}[H]
  \centering
  \caption{The $100\ \Ohm$ series element against record length, on one shared mesh: cell $0.400$ mm
  ($n_{\mathrm{sub}} = 4$), $n_{\mathrm{strip}} = 8$, $\eps_r = 4.4$, $h = 1.60$ mm, six matched-layer cells on
  every face, the cut at $15.20$ mm, the drive asked at $5.827786$ GHz. The loaded column is the reference-given
  reading, and the twin column is the same run with the load removed, on the identical mesh and at the identical
  step count. The $300$-step loaded and twin entries are the rows of Table~\ref{tab:loads-ratio}, and the
  $2400$-step entry is the $100\ \Ohm$ point of Table~\ref{tab:loads-series}. A dash marks a twin that was not
  run at that length.}
  \label{tab:loads-length}
  \begin{tabular}{@{}rrr@{}}
    \toprule
    $n_{\mathrm{steps}}$ & loaded, reference given & load-free twin \\
    \midrule
     300 & $0.613960$ & $0.541741$ \\
     600 & $0.517615$ & --- \\
    1200 & $0.543970$ & --- \\
    2400 & $0.549464$ & $0.981263$ \\
    4800 & $0.555663$ & $0.987637$ \\
    9600 & $0.556575$ & $0.989576$ \\
    \bottomrule
  \end{tabular}
\end{table}

The same readings can be checked against the physics the element is supposed to show. With the measured
$Z_c = 50.547 + 4.025j\ \Ohm$ of Table~\ref{tab:loads-series}, a $100\ \Ohm$ series element gives
$|2Z_c/(2Z_c + 100)| = 0.503908$, which is exactly the predicted column printed in that table; the measured
deviation from it is $+9.040$ per cent at $2400$ steps, $+10.272$ per cent at $4800$ steps and $+10.453$ per cent
at $9600$ steps, and the limit near $0.5567$ corresponds to about $+10.5$ per cent. The deviation therefore
converges \emph{away} from the ideal, and the $9.040$ per cent printed in Table~\ref{tab:loads-series} is an
understatement of the model's bias by roughly $1.4$ percentage points. That column is left as it stands.

The drive frequency is not the explanation: the metrics row that reports the actual drive frequency reads
$5.8278$ GHz at $300$, $600$, $4800$ and $9600$ steps alike, so the DFT-bin snap did not move the excitation, and
record length is the only variable that moved.

The reference path is an identity between the first and third rows, and the table is quoted as a check rather than
as a measurement. Dividing the unreferenced reading by the twin's own reading reproduces the referenced one,
\[
  \frac{0.327195}{0.541741} = 0.603970,
\]
which reproduces $0.613960$ to within $1.627$ per cent. That $1.627$ per cent is a residual difference between the
fit and the estimator, not agreement: the division is exact only if the twin supplies the line's delay and impedance
without error, and what is left over is the fit's own scatter carried through the ratio. As a residual it is the
check the table exists for -- a stale reference, a different pair of stations or a different reference impedance
breaks the identity by far more than a per cent -- but it is not a statement that the two numbers agree, and it is
not quoted as one. What the check certifies is that the three readings come from one record and one mesh, worked
three ways; both of the numbers it divides are themselves unsettled in record length, so an identity that holds
between two moving readings is a statement about the arithmetic and not about the line.

The third row is what an unreferenced reading of a loaded record looks like: $0.327195$, with no normalisation at
all, so it carries whatever the record itself contributes. On this record that contribution is large -- the
load-free twin reads $0.541741$ at $300$ steps -- and Table~\ref{tab:loads-length} shows that most of that
$0.541741$ is record length rather than line loss, since the same twin reaches $0.989576$ at $9600$ steps. Dividing
by the twin is still what the arithmetic requires before the reading can be called a transmission coefficient, but
on a $300$-step record it divides out a number that is itself unsettled, which is why the result above is quoted as
a consistency check.

\subsection{The shunt branch, and how far it gets}
\label{sec:loads-shunt}

A \texttt{shunt\_c} load is not a lumped capacitor and is not claimed to be one. It is a permittivity injection
on a fixed patch of $E_z$ edges at the load column, and the capacitance a user commands is realised only through
that injection. Table~\ref{tab:loads-shunt-cap} sweeps the command over a factor of forty on the same $0.400$ mm
mesh and reports what comes back, together with the depth of the injection expressed as the injected per-edge
capacitance divided by the cell's own.

\begin{table}[H]
  \centering
  \caption{Commanded against recovered capacitance for the \texttt{shunt\_c} element on the $0.400$ mm mesh, and
  the injection depth the command corresponds to. The cell's own capacitance is
  $\eps_0\eps_r\,\dd x = 1.5583\times10^{-14}$ F at $\eps_r = 4.4$, and the depth row is the injected per-edge
  capacitance divided by it, so it is dimensionless. The recovery is exact in bookkeeping and increasingly
  overstated in field: at $1$ pF the recovered capacitance exceeds the commanded one by a quarter. The
  parallel-plate harness of Section~\ref{sec:testing}, where the same element recovers $+0.31$ per cent at an
  injection depth of $4.58$, is the positive control this table is read against.}
  \label{tab:loads-shunt-cap}
  \begin{tabular}{@{}lrrrrrr@{}}
    \toprule
    Commanded $C$ [pF] & 0.05 & 0.1 & 0.2 & 0.5 & 1.0 & 2.0 \\
    \midrule
    Recovered $C$ [pF] & 0.049771 & 0.100214 & 0.203467 & 0.540617 & 1.250664 & 3.716293 \\
    Ratio              & 0.9954   & 1.0021   & 1.0173   & 1.0812   & 1.2507   & 1.8581   \\
    Injection depth    & 1.60     & 3.21     & 6.42     & 16.04    & 32.09    & 64.17    \\
    \bottomrule
  \end{tabular}
\end{table}

The edge bookkeeping is exact, and was checked against the arrays rather than assumed. The load column carries
$32$ $E_z$ edges -- eight ground-to-strip columns of four series edges each -- the per-edge permittivity is raised
by $C\,n_{\mathrm{sub}}/n_{\mathrm{strip}} = C/2$, and the series-parallel combination of eight columns of four
series edges returns exactly $C$. The deviation in the table is therefore not a counting error. Nor is it a fixed
parasitic: the recovered value exceeds the commanded one by $0.0035$ pF at $0.2$ pF and by $1.716$ pF at $2.0$ pF,
where a fixed offset would give the same number twice. Nor is it a multiplicative factor, since the ratio spans
$0.9954$ to $1.8581$.

What a user actually sees is the loading itself, in Table~\ref{tab:loads-shunt-load}.

\begin{table}[H]
  \centering
  \caption{The loading the \texttt{shunt\_c} element produces on the reference board of
  Table~\ref{tab:loads-series}: measured transmission coefficient against the commanded capacitance, same board,
  same step count, same load-free reference. The shallowest command moves the transmission by less than the
  resolution of the measurement, and the deepest removes three quarters of it while rotating the phase by
  $73^\circ$.}
  \label{tab:loads-shunt-load}
  \begin{tabular}{@{}lrrrrrr@{}}
    \toprule
    Commanded $C$ [pF] & 0.05 & 0.1 & 0.2 & 0.5 & 1.0 & 2.0 \\
    \midrule
    $|s_{21}|$             & 1.000292 & 0.998486 & 0.988121 & 0.901756 & 0.656036 & 0.276753 \\
    $\arg s_{21}$ [$^\circ$] & $-2.6513$ & $-5.3342$ & $-10.7653$ & $-26.9490$ & $-49.7584$ & $-73.4385$ \\
    \bottomrule
  \end{tabular}
\end{table}

The loading is monotone and smooth, and at the shallowest command the element is indistinguishable from nothing:
$1.000292$ is three hundredths of a per cent above unity, which is the resolution of the measurement rather
than a gain, and it is the only point in the table that reads above unity.

The cause of the overstatement is identifiable; the mechanism is not closed, and this paragraph is the place
where that distinction is kept. At $1$ pF the branch overstates the commanded capacitance by $25$ per cent
(ratio $1.2507$) and does not recover it, and a sweep on one mesh cannot say why. On a fixed mesh the commanded
capacitance, the injection depth and the perturbed phase advance per cell are all monotone functions of one
another, so any of the three explains the column equally well; a curve fitted to this table would be a
re-parameterisation of the sweep and not evidence about its cause. Refinement does not rescue the explanation
either: across meshes the behaviour is not consistent with a phase-advance reading at all. What does survive as
evidence is the positive control named in the caption of Table~\ref{tab:loads-shunt-cap}, where the same
injection recovers $+0.31$ per cent at a depth within a factor of two of what this mesh gives at a matched
depth. The structural difference between the two is the shape of the injection: in the harness it spans the whole
cross-section and therefore changes the medium, whereas on the microstrip it spans only the strip footprint,
inside the non-uniform fringing field of an $\eps_r = 4.4$ substrate, and perturbs the field distribution as well
as the stored energy. No single scalar captures a difference in kind.

The conclusion is therefore an envelope and not a validation. On this mesh the \texttt{shunt\_c} element
recovers the commanded capacitance to a few tenths of a per cent only while the injection is shallow -- it is
$0.46$ per cent low at $0.05$ pF and $0.21$ per cent high at $0.1$ pF, already $1.73$ per cent high at $0.2$ pF,
which is $6.42$ cell capacitances per edge -- and by $1$ pF it overstates by a quarter. The cause is named and
the mechanism is not closed, so \emph{the shunt branch is not claimed to be validated}. It is recorded here in the
same way as the other open residual of this work, rather than presented as a success because the shallow points
look good.

The rule that follows is a rule for use. Command a \texttt{shunt\_c} element at or below a few times the cell's
own capacitance and it does what it says to a fraction of a per cent. Command a larger one on a mesh this coarse
and what is being set is the strength of an injection, not a capacitance: the element will load the line, the
loading will be smooth and monotone as Table~\ref{tab:loads-shunt-load} shows, and the number a user asked for
will not be the number the line meets.
